\begin{problem}
    Para arrendar botes en una empresa se debe enviar una solicitud. En la siguiente figura se representa una solicitud que está en estado de “Aprobada” y le falta pasar por los últimos tres estados.

    \begin{center}
    (imagen)
    \end{center}

    Para calcular el progreso de las últimas \num{20} solicitudes, la empresa calcula el porcentaje de ellas que está o ya ha pasado el estado “Aprobada”.

    En la siguiente tabla se presentan las últimas \num{20} solicitudes:

    \[
    \begin{array}{c|c|c|c}
    \text{Solicitud} & \text{Estado} & \text{Solicitud} & \text{Estado} \\
    \hline
    1 & \text{Enviada} & 11 & \text{En curso} \\
    2 & \text{En curso} & 12 & \text{En curso} \\
    3 & \text{En curso} & 13 & \text{En curso} \\
    4 & \text{Finalizada} & 14 & \text{Enviada} \\
    5 & \text{Finalizada} & 15 & \text{Enviada} \\
    6 & \text{Pagada} & 16 & \text{Pagada} \\
    7 & \text{Enviada} & 17 & \text{Finalizada} \\
    8 & \text{Aprobada} & 18 & \text{Aprobada} \\
    9 & \text{Enviada} & 19 & \text{Pagada} \\
    10 & \text{Pagada} & 20 & \text{Pagada}
    \end{array}
    \]

    ¿Cuál es el porcentaje de las solicitudes que está o ya ha pasado por el estado “Aprobada”?

    \begin{enumerate}[label=\Alph*.]
        \item \SI{13}{\%}
        \item \SI{15}{\%}
        \item \SI{65}{\%}
        \item \SI{75}{\%}
    \end{enumerate}
\end{problem}
